\section{Related literature}
\label{sec:literature}

\paragraph{The FRB/US lineage.} FRB/US descends from a half-century of structural modelling at the Federal Reserve Board. Its predecessor, the MIT--Penn--SSRC (MPS) model, served as the Board's workhorse from around 1970; \citet{brayton1997evolution} recount how accumulating dissatisfaction with the MPS treatment of expectations and dynamics --- essentially unrestricted distributed lags, exposed to the Lucas critique --- motivated the 1996 redesign. FRB/US's two signature choices date from that redesign. First, dynamics are derived from explicit adjustment-cost optimisation: the polynomial-adjustment-cost framework of \citet{tinsley1993, tinsley2002} generalises quadratic adjustment costs to higher-order frictions, delivering estimated error-correction decision rules that remain interpretable as optimising behaviour (Section~\ref{sec:pac}). Second, expectations are explicit and switchable: every behavioural equation separates its expectational inputs into named variables that can be generated either by small estimated VARs or model-consistently, so the same estimated structure supports both boundedly rational and rational-expectations analysis \citep{brayton1997role}. \citet{brayton2014note} is the standard modern summary and the source of the published impulse responses used for comparison in Section~\ref{sec:results}.

\paragraph{Policy analysis with FRB/US.} The model has been the vehicle for a long line of influential policy research. \citet{reifschneider2000three} used stochastic simulations of FRB/US to show that the zero lower bound becomes quantitatively important for inflation targets around or below 2 per cent --- under a Taylor rule with a 2 per cent target the funds rate is constrained roughly a tenth of the time --- and proposed the ``lower-for-longer'' compensating rules that later underpinned makeup-strategy analysis. \citet{laubach2003measuring}, estimating the natural real rate jointly with potential output, supplied the $r^{*}$ concept that enters FRB/US policy rules as \texttt{rstar} and conditions its low-rate scenario work. \citet{reifschneider2016gauging} used FRB/US to ask whether asset purchases and forward guidance could substitute for constrained rate cuts in a future recession, and \citet{bernanke2019strategies} ran large stochastic-simulation comparisons of Taylor, Reifschneider--Williams and temporary price-level-targeting rules in FRB/US under both expectations options, finding makeup strategies robust to imperfect credibility. The model remained central to the Federal Reserve's 2019--2020 framework review --- \citet{hebden2020makeup} stress-test average-inflation-targeting and price-level-targeting strategies in FRB/US across its full expectations menu, finding that under VAR expectations makeup strategies provide no early stabilisation and later overheat the economy --- and to its aftermath: \citet{brayton2022linver} distil FRB/US into a linearised version (LINVER) for large-scale stochastic ELB simulations, and \citet{kiley2025anchoring} uses FRB/US to quantify how incorporating long-run inflation expectations into policy rules affects anchoring, the most recent Board research use of the model at the time of writing. On the fiscal side, FRB/US was one of the models in the multi-country fiscal-stimulus comparison of \citet{coenen2012effects}, which reports first-year U.S.\ government-consumption multipliers of 0.7--1.0 without monetary accommodation (FRB/US near the bottom of that range, rising to roughly 1.1--1.2 with two years of accommodation) and labour-tax-cut multipliers of only 0.2--0.4; \citet{cashin2018fiscal} construct a FRB/US-based ``fiscal effect'' measure of discretionary policy around the Great Recession, with personal-tax MPCs of a few tenths. These published magnitudes are the yardsticks for the experiments of Section~\ref{sec:results}, alongside the survey of multiplier evidence in \citet{whalen2015fiscal}.

\paragraph{Open-source macromodelling.} FRB/US is unusual among central-bank workhorse models in being fully public: equations, coefficients, database, documentation and simulation code are all in the public domain \citep{fedfrbusabout}. The Board's own \texttt{pyfrbus} package \citep{fedfrbuspython} moved the platform from EViews to Python in 2021, substantially widening access; but, as Section~\ref{sec:implementation} details, it is pinned to a scientific-Python stack of that era. This paper's contribution sits in the small but growing space of independently engineered, continuously tested open implementations of official policy models --- the macro counterpart to open tax--benefit microsimulation platforms such as PolicyEngine's, with which it shares tooling conventions (public artefacts as the single source of truth, validation gates in continuous integration, semantic versioning of data vintages). The intended payoff is composability: microsimulation-based revenue and distributional estimates on one side, and model-consistent macroeconomic feedbacks from a published, Fed-maintained structure on the other.
